\section*{Hoofdstuk \ref{ch:b1:plant-water-minerals} --- Water- en mineralenvoorziening van planten}

\begin{solution}{exo:b1:plant-water-minerals:1}
\omterm{def:b1:plant-water-minerals:pathways}{Apoplast}: de aaneengesloten wanden en holten, doorkruist zonder een
membraan te passeren. \omterm{def:b1:plant-water-minerals:pathways}{Symplast}: het verbonden \omterm{def:b1:cell-unit-of-life:cell}{cytoplasma} van de \omterm{def:b1:cell-unit-of-life:cell}{cellen}, dat
door één membraan wordt betreden en via plasmodesmata wordt voortgezet. De
\omterm{def:b1:plant-water-minerals:pathways}{band van Caspary} verspert de \omterm{def:b1:plant-water-minerals:pathways}{apoplast} bij de \omterm{prop:b1:flowering-plant-organization:stemroot}{endodermis}, zodat alles door
een membraan moet en er niets uit het \omterm{def:b1:flowering-plant-organization:tissues}{xyleem} terug kan lekken.
\end{solution}

\begin{solution}{exo:b1:plant-water-minerals:2}
N (\omterm{def:b1:proteins:peptide}{eiwitten}, nucleïnezuren), P (\omterm{def:b1:nucleic-acids:nucleotide}{nucleotiden}, membranen), K (osmotisch
middel, huidmondjes), Ca (wanden, signalering), Mg (\omterm{def:b1:photosynthesis:pigments}{chlorofyl}), S (de
\omterm{def:b1:proteins:aminoacid}{aminozuren} cysteïne en methionine).
\end{solution}

\begin{solution}{exo:b1:plant-water-minerals:3}
Bodem \qty{-0.2}{}, wortelxyleem \qty{-0.6}{}, bladcel \qty{-1.3}{MPa}: van
de bodem de \omterm{def:b1:flowering-plant-organization:organs}{wortel} in, van het wortelxyleem omhoog naar het \omterm{def:b1:flowering-plant-organization:organs}{blad} --- altijd
naar de negatievere waarde toe.
\end{solution}

\begin{solution}{exo:b1:plant-water-minerals:4}
Stikstof is beweeglijk in het floëem: de plant onttrekt haar aan oude
\omterm{def:b1:flowering-plant-organization:organs}{bladeren} om de jonge te voeden, zodat de oude \omterm{def:b1:flowering-plant-organization:organs}{bladeren} het eerst vergelen.
IJzer wordt niet opnieuw gemobiliseerd: de jonge \omterm{def:b1:flowering-plant-organization:organs}{bladeren}, zonder ijzer
gebouwd, zijn degene die vergelen.
\end{solution}

\begin{solution}{exo:b1:plant-water-minerals:5}
\omterm{def:b1:cell-unit-of-life:cell}{Cel} $\Psi = -0.9 + 0.5 = \qty{-0.4}{MPa}$: bodem ($-0.2$) $>$ \omterm{def:b1:cell-unit-of-life:cell}{cel}
($-0.4$) $>$ \omterm{def:b1:flowering-plant-organization:tissues}{xyleem} ($-0.6$), dus stroomt water bodem $\to$ \omterm{def:b1:cell-unit-of-life:cell}{cel} $\to$
\omterm{def:b1:flowering-plant-organization:tissues}{xyleem}. De opname stopt wanneer de \omterm{def:b1:cell-unit-of-life:cell}{cel} \qty{-0.2}{MPa} bereikt, dat wil
zeggen $\Psi_p = -0.2 + 0.9 = \qty{0.7}{MPa}$.
\end{solution}

\begin{solution}{exo:b1:plant-water-minerals:6}
$\Psi = 137\ln(\text{RV})$: \qty{-14}{MPa}, \qty{-95}{MPa},
\qty{-315}{MPa}. Zelfs vochtige lucht is in deze zin tien keer droger dan
een bodem op het verwelkingspunt: een \omterm{def:b1:flowering-plant-organization:organs}{blad} verliest altijd water aan de
lucht; de enige vraag is hoe snel.
\end{solution}

\begin{solution}{exo:b1:plant-water-minerals:7}
$E_K = 25.7\ln(0.1/100) = \qty{-178}{mV}$: bij \qty{-150}{mV} is de \omterm{def:b1:cell-unit-of-life:cell}{cel}
iets minder negatief dan nodig, dus een duizendvoudige ophoping vraagt iets
meer dan \omterm{def:b1:membranes-transport:transporters}{kanalen} --- maar een verhouding van 100 zou in evenwicht zijn:
\omterm{def:b1:membranes-transport:transporters}{kanalen} doen bijna alles. Nitraat: $E = -25.7\ln(0.5/5)
= \qty{+59}{mV}$; de \omterm{def:b1:cell-unit-of-life:cell}{cel} staat op \qty{-150}{}, \qty{209}{mV} daarvandaan
en in de verkeerde richting: een anion moet naar binnen worden gepompt,
door protonsymport.
\end{solution}

\begin{solution}{exo:b1:plant-water-minerals:8}
's Nachts, zonder transpiratie, trekt het laden van ionen in het \omterm{def:b1:flowering-plant-organization:tissues}{xyleem}
water aan en duwt een \omterm{prop:b1:plant-water-minerals:rootpressure}{worteldruk} van enkele tienden van een megapascal het
sap hoogstens enkele meters omhoog --- genoeg om de toppen van grassprieten
te bereiken en er druppels uit te persen, niet om de top van een boom te
bereiken. Zodra de huidmondjes opengaan, zet de transpiratie het \omterm{def:b1:flowering-plant-organization:tissues}{xyleem}
onder spanning en verdwijnt de \omterm{prop:b1:plant-water-minerals:rootpressure}{worteldruk}.
\end{solution}

\begin{solution}{exo:b1:plant-water-minerals:9}
Zonder molybdeen werkt \omterm{prop:b1:plant-water-minerals:nitrate}{nitraatreductase} niet: op nitraat kan de plant geen
\omterm{def:b1:proteins:aminoacid}{aminozuren} maken en vertoont zij een stikstofgebrek hoewel er nitraat in
overvloed is. Op ammonium groeit zij normaal, omdat ammonium zonder
reductie in de \omterm{def:b1:proteins:aminoacid}{aminozuren} terechtkomt. Het verschil legt het element in één
\omterm{def:b1:enzymes:enzyme}{enzym} vast.
\end{solution}

\begin{solution}{exo:b1:plant-water-minerals:10}
\qty{30}{mg} N $= \qty{2.14}{mmol}$; $\times 8 = \qty{17}{mmol}$
elektronen, \qty{8.6}{mmol} NADPH-equivalenten. \omterm{def:b1:flowering-plant-organization:organs}{Blad}: $0.05\times
4\times 10^{-6}\times 43\,200 = \qty{8.6}{mmol}$ $\mathrm{CO_2}$,
\qty{35}{mmol} elektronen: de nitraatreductie vraagt op deze telling half
zoveel elektronen als de koolstofvastlegging --- een groot deel, en daarom
doen planten het in het licht en zijn snelgroeiende \omterm{def:b1:flowering-plant-organization:organs}{bladeren} grootverbruikers
van reductiemiddel.
\end{solution}

\begin{solution}{exo:b1:plant-water-minerals:11}
\qty{2.14}{mmol} N $= \qty{1.07}{mmol}$ $\mathrm{N_2}$: \omterm{def:b1:water-small-molecules:atp}{ATP}
$16\times 1.07 = \qty{17}{mmol}$, dat wil zeggen \qty{0.57}{mmol} glucose,
\qty{0.10}{g}; plus 8 elektronen per $\mathrm{N_2}$ (nog eens
\qty{0.05}{g} glucose). De regel van \qty{6}{g/g} geeft \qty{0.18}{g} per
dag: nitrogenase zelf is er het meeste van, en de rest is het bouwen en
onderhouden van de knolletjes, hun leghemoglobine, en de eigen ademhaling
van de bacteriën om de zuurstof laag te houden.
\end{solution}

\begin{solution}{exo:b1:plant-water-minerals:12}
De \omterm{def:b1:flowering-plant-organization:organs}{wortel} neemt zelf het water en de beweeglijke ionen (nitraat, kalium)
op die het bodemwater tot bij zijn haren brengt, en besteedt alleen \omterm{def:b1:water-small-molecules:atp}{ATP} aan
de \omterm{prop:b1:plant-water-minerals:uptake}{protonpomp}. Voor fosfaat, dat niet naar hem toe komt, betaalt hij een
schimmel in suiker om zijn bereik te verlengen; voor stikstof in een bodem
zonder nitraat betaalt hij bacteriën in suiker om die te binden. Het
honorarium is een tiende of meer van zijn fotosynthese, en de plant stopt
met betalen zodra de dienst niet nodig is: op fosfaatrijke bodem worden de
\omterm{def:b1:plant-water-minerals:mycorrhiza}{mycorrhiza}'s afgebouwd, op nitraatrijke bodem maken vlinderbloemigen minder
knolletjes --- aannemers blijven alleen zolang zij goedkoper zijn dan het
werk zelf doen.
\end{solution}

\begin{solution}{pb:b1:plant-water-minerals:1}
\textbf{1.} $-0.9 + 0.5 = \qty{-0.4}{MPa}$.
\textbf{2.} $-0.2 > -0.4 > -0.6 > -1.3$: elke stap lager, het water beweegt
bodem $\to$ schors $\to$ \omterm{def:b1:flowering-plant-organization:tissues}{xyleem} $\to$ \omterm{def:b1:flowering-plant-organization:organs}{blad}.
\textbf{3.} $137\ln 0.5 = \qty{-95}{MPa}$.
\textbf{4.} Bodem naar \omterm{def:b1:flowering-plant-organization:organs}{blad} \qty{1.1}{MPa}; \omterm{def:b1:flowering-plant-organization:organs}{blad} naar lucht \qty{94}{MPa}:
\qty{99}{\%} van het geheel is de laatste stap.
\textbf{5.} Een bodem op $-0.9$ ligt onder de $-0.4$ van de schorscellen:
er verlaat water de \omterm{def:b1:cell-unit-of-life:cell}{cellen} tot zij op $-0.9$ komen, dat wil zeggen tot hun
\omterm{def:b1:eukaryotic-cell:plantcell}{turgor} tot nul daalt; de opname stopt tot de \omterm{def:b1:cell-unit-of-life:cell}{cellen} hun $\Psi_s$ verlagen
door opgeloste stoffen op te hopen.
\textbf{6.} De bodem ligt nu onder de $-1.3$ van de \omterm{def:b1:flowering-plant-organization:organs}{bladeren}: er stroomt
water uit de plant; de \omterm{def:b1:flowering-plant-organization:organs}{bladeren} verliezen hun \omterm{def:b1:eukaryotic-cell:plantcell}{turgor} en verwelken. Water
geven tilt de $\Psi$ van de bodem boven die van de \omterm{def:b1:flowering-plant-organization:organs}{wortel} en de stroom
hervat zich; de \omterm{def:b1:flowering-plant-organization:organs}{bladeren}, waarvan de wanden ongeschonden waren, vullen zich
weer.
\textbf{7.} $2\times 10^{-7}\times 0.2\times(-0.2 - (-0.6)) =
\qty{1.6e-8}{m^3/s}$.
\textbf{8.} $1.6\times 10^{-8}\times 86\,400 = \qty{1.4e-3}{m^3} =
\qty{1.4}{L}$ per dag.
\textbf{9.} \qty{1400}{g} over \qty{0.5}{m^2} en \qty{12}{h} daglicht (het
meeste ervan): ongeveer \qty{230}{g} per vierkante meter per uur.
\textbf{10.} $230/18 = \qty{12.8}{mol}$ per vierkante meter per uur, \qty{3.5}{mmol} per vierkante meter per seconde --- duizend watermoleculen voor elke $\mathrm{CO_2}$ die bij \qty{4}{\micro mol} wordt vastgelegd.
\textbf{11.} Verschil $-0.2 - (-0.4) = 0.2$: de flux is gehalveerd,
\qty{0.7}{L} per dag. Het sluiten van de huidmondjes heeft het verlies
gesnoeid, de spanning is verslapt, en de hele ladder is gestegen: de plant
ruilt koolstofwinst voor water wanneer de lucht het droogst is.
\textbf{12.} Water beweegt langs zijn eigen potentiaalgradiënt, die door de
verdamping van het \omterm{def:b1:flowering-plant-organization:organs}{blad} wordt gemaakt, zodat de \omterm{def:b1:flowering-plant-organization:organs}{wortel} lijdelijk is; nitraat
moet tegen zowel een concentratie- als een elektrische gradiënt in worden
geconcentreerd, wat alleen de koppeling aan de \omterm{prop:b1:plant-water-minerals:uptake}{protonpomp} kan betalen.
\textbf{13.} $\qty{1.4}{L}\times\qty{2}{mmol/L} = \qty{2.8}{mmol}$,
\qty{39}{mg} stikstof.
\textbf{14.} Ja, met een derde over; het overschot wordt als nitraat in de
\omterm{def:b1:eukaryotic-cell:plantcell}{vacuolen} opgeslagen of bij de \omterm{prop:b1:flowering-plant-organization:stemroot}{endodermis} buitengehouden.
\textbf{15.} $\Delta G = RT\ln(5/2) + zFV = 2.48\times 0.92 + (-1)(96.5)(-0.15)
= 2.3 + 14.5 = \qty{16.8}{kJ/mol}$.
\textbf{16.} Per proton: $RT\ln(10^{-7.5}/10^{-5.5}) + F(-0.15) =
-11.4 - 14.5 = \qty{-25.9}{kJ}$; twee protonen: \qty{-52}{kJ}, drie keer
wat er nodig is.
\textbf{17.} $\qty{2.14}{mmol}\times 8 = \qty{17}{mmol}$ elektronen.
\textbf{18.} $0.5\times 4\times 10^{-6}\times 43\,200 = \qty{86}{mmol}$
$\mathrm{CO_2}$, \qty{346}{mmol} elektronen: \qty{5}{\%} naar
nitraat.
\textbf{19.} Het zou de 8 elektronen per stikstofatoom besparen; maar het
opnemen van een kation en het uitstoten van een proton per stuk verzuurt de
bodem rond de \omterm{def:b1:flowering-plant-organization:organs}{wortel} en de \omterm{def:b1:cell-unit-of-life:cell}{cel} moet \omterm{def:b1:water-small-molecules:ph}{zuur} uitvoeren, en ammonium is giftig
als het zich ophoopt, zodat het meteen in de \omterm{def:b1:flowering-plant-organization:organs}{wortel} moet worden
geassimileerd.
\textbf{20.} \qty{1.07}{mmol} $\mathrm{N_2}$ $\times 16 =
\qty{17}{mmol}$ \omterm{def:b1:water-small-molecules:atp}{ATP}.
\textbf{21.} $6\times 0.03 = \qty{0.18}{g}$ suiker, \qty{6}{\%} van
\qty{3}{g}. Met \qty{2}{mmol/L} nitraat beschikbaar voor de kosten van de
reductie ervan (\qty{5}{\%} van de elektronen) is binding een slechte ruil
en maakt een vlinderbloemige weinig knolletjes; zonder nitraat is het de
enige weg, en is \qty{6}{\%} van de productie goedkoop voor het volledige
\omterm{def:b1:proteins:peptide}{eiwit} van de plant.
\textbf{22.} $1.4\times 2\times 10^{-6} = \qty{2.8}{\micro mol}$,
\qty{0.09}{mg} fosfor --- een veertigste van de behoefte; de plant moet
steunen op de diffusie naar haar \omterm{def:b1:flowering-plant-organization:organs}{wortelharen}, die hun omgeving spoedig
uitputten, en vooral op haar mycorrhizaschimmel.
\textbf{23.} De band verspert de \omterm{def:b1:plant-water-minerals:pathways}{apoplast}, zodat nitraat het \omterm{def:b1:flowering-plant-organization:tissues}{xyleem} alleen
kan bereiken door membranen van de \omterm{prop:b1:flowering-plant-organization:stemroot}{endodermis} die het naar binnen pompen; en
hij verhindert dat het geconcentreerde xyleemsap langs de wanden terug
weglekt, zodat de gradiënt behouden blijft.
\textbf{24.} Zijn xyleemsap zou op de bodemoplossing lijken --- verdund,
ongeselecteerd, met natrium en andere ongewenste ionen --- en kon niet
geconcentreerd worden gehouden; in zoute bodem zou natrium vrij de \omterm{def:b1:flowering-plant-organization:organs}{bladeren}
bereiken en zou de plant vergiftigd raken, waar een normale plant het bij de
\omterm{prop:b1:flowering-plant-organization:stemroot}{endodermis} buitenhoudt.
\textbf{25.} Ongeveer \qty{1.4}{L} water per dag door \qty{0.2}{m^2}
\omterm{def:b1:flowering-plant-organization:organs}{wortel}, dat \qty{39}{mg} stikstof aanvoert tegen een behoefte van
\qty{30}{}, waarvan de reductie ongeveer \qty{5}{\%} van de
fotosynthetische elektronen van het \omterm{def:b1:flowering-plant-organization:organs}{blad} vraagt.
\end{solution}
